  \section{Lacunes de recherche et opportunités}\label{sec:lacunes-recherche}
    Les sections précédentes ont détaillé les coûts, les impacts environnementaux et les divers effets du stationnement sur la mobilité des personnes et la forme urbaine. L'offre de stationnement sur rue a jusqu'à présent reçu plus d'attention \parencite{bourdeau_methodologie_2014,open_mobility_foundation_curb_2024-1,eros_curblr_2025} puisqu'il constitue la majorité l'offre de stationnement dans les lieux où il y a une forte demande et les plus hauts taux d'occupation \parencite{morency_stationnement_2017,aignel_elaboration_2015}, et que c'est la responsabilité directe des gouvernements municipaux de gérer ce parc.\par

  
    Cependant, plusieurs études montrent que le stationnement hors rue peut avoir une forte influence sur la mobilité des personnes même en présence d'alternatives viables \parencite{guo_does_2013,currans_households_2023,weinberger_residential_2009}. Malgré cela, il existe peu d'information détaillée sur l'offre de stationnement sur l'ensemble du territoire québécois. Bien qu'un dénombrement montréalais existe pour les stationnements non résidentiels \parencite{consortium_cima_-_daniel_arbour_et_associes_inventaire_1998,aignel_elaboration_2015}, cette information est récoltée de manière manuelle et n'inclut pas le stationnement résidentiel. D'autre part, \textcite{chester_inventorying_2022} met à disposition un fichier Excel des règlements qui intègre un format de données\parencite{li_sfba_2022} intégrant les unités, il n'intègre pas toutes les variations constatées pour l'étude de cas et la variation temporelle des règlements n'est pas incluse. Même si certaines études ont utilisé les minimums de stationnements pour déduire une capacité de stationnement \parencite{chester_parking_2015,chester_inventorying_2022,hoehne_valley_2019}, les trois articles montrent des niveaux de complexité différents et ne sont pas explicites sur comment la variation temporelle des règlements est intégrée. D'autre part, bien que des efforts aient été mis en place pour numériser d'autres parties de la réglementation d'urbanisme \parencite{xu_national_2023}, la compilation d'information sur la réglementation de stationnement se résume actuellement à l'abrogation des minimums de stationnement \parencite{parking_reform_network_parking_2025,bronin_how_2023}. L'absence de données sur ces règlements rend difficiles la comparaison de différents régimes réglementaires et l'illustration des coûts d'opportunités et des potentielles incohérences de ces politiques.\par
    La mesure de phénomène et de quantités est centrale à la démarche scientifique, au génie et à la gestion en général \parencite{kelvin_popular_1891,taylor_principles_1911,dolence_using_2015}. Lord Kelvin le statuait comme suit:
    \begin{quote}
        [...] Une première étape essentielle dans l'apprentissage de tout sujet est de trouver des mécanismes de calcul numérique et des méthodes pratiques de mesure d'une quantité reliée au sujet. Je dis souvent que quand on mesure ce dont on parle et qu'on peut l'exprimer en chiffres, on en connaît quelque chose; mais quand on ne peut le mesurer, quand on ne peut l'exprimer en chiffres, notre connaissance est d'une nature maigre et insatisfaisante
    \end{quote}
    Le stationnement a largement échappé à ce précepte. Cela est dû à la difficulté de la collecte des données reliées à la réglementation \parencite{axelrod_automating_2023}, mais aussi à la collation de données fiables pour le parc existant \parencite{morency_stationnement_2017,heran_consommation_2008}. Plusieurs sources \parencite{morency_stationnement_2017,xu_national_2023,bronin_how_2023,currans_households_2023} indiquent que la mise en place d'une plateforme de gestion et de communication d'information reliée à la réglementation, mais aussi à l'offre actuelle est un élément indispensable pour comprendre les usages et arbitrer leur partage dans l'espace urbain.\par
    D'autre part, même s'il existe plusieurs formats pour représenter le stationnement\parencite{eros_curblr_2025,osm_contributors_amenity_2025,organisation_internationale_de_normalisation_intelligent_2023}, tous font l'hypothèse qu'il existe un estimateur précis pour la capacité de stationnement. Or plusieurs des méthodes impliquent le respect d'hypothèses ou un degré d'erreur. De plus, aucune ne prend en compte la possibilité de plusieurs estimateurs différents donnant des valeurs différentes pour un même lieu. \par
    Finalement, bien qu'il existe plusieurs outils en source ouverte \parencite{osm_contributors_openstreetmap_2025,eros_curblr_2025,open_mobility_foundation_curb_2024-1, seidel_methodenbericht_2021} ou propriétaires \parencite{arcadis_curbiq_2025,populus_curb_2025} pour gérer des dénombrements de stationnement, aucun n'a été recensé pour faire le recensement du stationnement hors rue sans requérir la modélisation manuelle de chaque aire de stationnement. Finalement, aucun outil ne permet de rapidement comprendre l'offre combinée du stationnement hors rue et sur rue (sans avoir besoin d'onéreux processus manuels). Lorsqu'ils existent, ces outils sont souvent utilisés pour donner un portrait agrégé de l'offre plutôt que ce à quoi ferait face un voyageur.\par
    Pour conclure, cette revue de littérature a révélé sept grandes lacunes dans les connaissances:
    \begin{enumerate}
        \item L'actuel dénombrement de stationnement hors rue pour la ville de Montréal requiert une collation manuelle et ne représente que le stationnement non résidentiel. D'autre part, il n'existe à la connaissance de l'auteur pas d'autre juridiction québécoise ayant estimé la capacité de stationnement. 
        \item Le schéma de données de \textcite{li_sfba_2022} est spécifique au contexte de son étude \parencite{chester_inventorying_2022} et n'intègre pas des éléments qui ont été constatés dans les règlements de la ville de Québec discutés au chapitre suivant
        \item Bien que des schémas existent pour la représentation de la capacité de stationnement, aucun ne prend en compte la possibilité que différentes méthodes d'estimation puissent être utilisées. D'autre part, lorsque la méthode réglementaire a été utilisée, les résultats ont été agrégés par aire de diffusion.
        \item La compilation des données sur le stationnement demeure onéreuse, manuelle et les méthodes de calcul sont souvent difficiles à mettre en place. Une fois ces inventaires complétés, aucune infrastructure pour le partage et surtout la correction et la mise à jour des données n'est mise en place
        \item Les études existantes montrent un lien significatif entre le stationnement et les comportements de mobilité désagrégés. Malgré cela, il n'existe à l'heure actuelle pas de modèle de causalité entre la provision de stationnements et le comportement des personnes comme cela a été le cas pour l'offre de transport en commun \parencite{deschaintres_measuring_2021}
        \item Bien qu'il existe d'exemple d'atlas ouverts permettant de comparer la réglementation d'urbanisme entre différentes juridictions, l'automatisation de la création de tels atlas demeure onéreuse. Certains chercheurs ont émis des résultats montrant que les grands modèles de langage puissent accélérer leur création, le processus demeure manuel du fait de la subtilité des codes d'urbanisme. D'autre part, les atlas existants ne donnent que des indications sur la présence ou l'absence de requis de stationnement sans en spécifier l'intensité s'ils existent. 
        \item Plusieurs chercheurs ont développé des méthodes d'identification de la capacité de stationnement à l'aide de méthodes d'\ac{AP}. Aucune étude n'a été recensée qui permette d'utiliser ces méthodes pour une région métropolitaine et d'en agréger des statistiques de manière ouvertes. Les études recensées se limitent à lister la performance des méthodes par rapport aux annotations manuelles.
    \end{enumerate}
    Les objectifs de recherche énumérés en introduction visent surtout à pallier les lacunes 1, 2, 3 et 4. Les autres points devront être poursuivis dans la recherche future.